\newcommand{\InlineFor}[2]{\State \textbf{for} #1 \textbf{do} #2}

\section{Pairing Operations as Polynomials}

Algorithm~\ref{alg:pairing} gives the pairing computation combining El Housni's R1CS optimizations~\cite{PR1CS22} with Novakovic-Eagen's residue witness technique~\cite{OPP24}. The 27th root of unity $W$ scales the Miller loop output $f$ to be a cubic residue (see~\cite{OPP24} §4.3); if $f$ is not already a cubic residue, multiplying by $W^w$, $w\in\{1,2\}$ corrects this.

Our contribution is to express each step of this algorithm as a single multi-operand polynomial product verified via the Schwartz-Zippel lemma, rather than as a sequence of individual binary multiplications as in~\cite{FXFM23}.

\begin{algorithm}
    \caption{Optimal Ate Multi-Pairing with Residue witness check}
    \label{alg:multipairing}
    \begin{algorithmic}[1]
    \Require Vectors $\vec{P} = (P_1, \dots, P_N) \in \G_1^N$, $\vec{Q} = (Q_1, \dots, Q_N) \in \G_2^N$, $c \in \Fe{12}, w \in \{0,1,2\}$
    \Statex $\text{Internal } W, W^2 \in \Fe{12}$ \Comment{$W$ is 27th root of unity}
    \Ensure $\prod_{k=1}^N a_{opt}(Q_k, P_k)$
    \State Write $s = 6t + 2 = \sum_{i=0}^{L-1} s_i 2^i$ where $s_i \in \{-1, 0, 1\}$ and $s_L = 1$
    \State $c^{-1} \leftarrow 1 / c$ \Comment{Inverse Residue witness}
    \State $f \leftarrow c^{-1}$ \Comment{Initialize with inverse residue witness}

    \InlineFor {$k = 1$ to $N$} {$T_k \leftarrow Q_k$}

    \Statex

    \For{$i = L - 2$ to $0$}
        \State $f \leftarrow f^2$ \label{pr:ml:sq}
        \State $f \leftarrow f \cdot \prod_{k=1}^N l_{T_k,T_k}(P_k)$ \label{pr:ml:ldbl}

        \InlineFor {$k = 1$ to $N$} {$T_k \leftarrow [2]T_k$}

        \If{$s_i = 1$}
            \State $f \leftarrow f \cdot c^{-1}$ \Comment{Residue witness}\label{pr:ml:rw1bit}
            \State $f \leftarrow f \cdot \prod_{k=1}^N l_{T_k,Q_k}(P_k)$ \label{pr:ml:ladd1bit}
            \InlineFor {$k = 1$ to $N$} {$T_k \leftarrow T_k + Q_k$}

        \ElsIf{$s_i = -1$}
            \State $f \leftarrow f \cdot c$ \Comment{Residue witness}\label{pr:ml:rw-1bit}
            \State $f \leftarrow f \cdot \prod_{k=1}^N l_{T_k,-Q_k}(P_k)$ \label{pr:ml:ladd-1bit}
            \InlineFor {$k = 1$ to $N$} {$T_k \leftarrow T_k - Q_k$}

        \EndIf
    \EndFor
    \Statex

    \InlineFor {$k = 1$ to $N$} {$Q_{1,k} \leftarrow \pi_p(Q_k)$, $Q_{2,k} \leftarrow \pi_p^2(Q_k)$}

    \State $f \leftarrow f \cdot \prod_{k=1}^N l_{T_k,Q_{1,k}}(P_k)$ \label{pr:ml:pi1}

    \InlineFor {$k = 1$ to $N$} {$T_k \leftarrow T_k + Q_{1,k}$}

    \State $f \leftarrow f \cdot \prod_{k=1}^N l_{T_k,-Q_{2,k}}(P_k)$ \label{pr:ml:pi2}

    \Statex
    \If{$w = 1$}
        \State $f \leftarrow f \cdot W$ \label{pr:ml:27r}
    \ElsIf{$w = 2$}
        \State $f \leftarrow f \cdot W^2$ \label{pr:ml:27r2}
    \EndIf

    \State $f \leftarrow f \cdot c^{-q} \cdot c^{q^2} \cdot c^{-q^3}$ \Comment{Uses Frobenius maps, Negative exponents use $c^{-1}$}\label{pr:ml:fe}

    \Statex

    \State \Return $f$
    \end{algorithmic}
\end{algorithm}

\subsection{Comparative Analysis~\label{sec:ops:comp-analysis}}

We compare the two schemes for a 3-pair Miller loop zero-bit step (the most common case in BN254's NAF\cite{naf} representation of $s = 6t+2$). Here $F \in \Fe{12}$ is the accumulator, $L_i \in \text{Sparse}_{01379}$ are the three doubling line functions, $R \in \Fe{12}$ is the updated accumulator, and $Q$ is the quotient polynomial. $\text{Sparse}_{01379}$ denotes the subspace of $\Fe{12}$ elements whose only nonzero coefficients are at degrees $0,1,3,7,9$---the form that BN254 line functions naturally take~\cite{PR1CS22}, reducing their evaluation cost from 12C to 4C.

Adapting~\cite{FXFM23} to lines \ref{pr:ml:sq}--\ref{pr:ml:ldbl} of Algorithm~\ref{alg:pairing} chains four binary multiplications, each producing an intermediate remainder that must be committed:

\begin{align*}
    F(x) \cdot F(x) &\stackrel{?}{=} Q_1(x) \cdot P_{12}(x) + T_1(x) \\
    T_1(x) \cdot L_1(x) &\stackrel{?}{=} Q_2(x) \cdot P_{12}(x) + T_2(x) \\
    T_2(x) \cdot L_2(x) &\stackrel{?}{=} Q_3(x) \cdot P_{12}(x) + T_3(x) \\
    T_3(x) \cdot L_3(x) &\stackrel{?}{=} Q(x) \cdot P_{12}(x) + R(x)
\end{align*}

Each of the five polynomial evaluations ($F$, $T_1$--$T_3$, $R$) costs 12C and each line costs 4C, with all four remainders requiring commitment: \textbf{76C and 48 $\Fq$ commitments}.

Our scheme collapses this into a single identity:\label{sec:op:comp:our}

$$F(x)^2 \cdot L_1(x) \cdot L_2(x) \cdot L_3(x) \stackrel{?}{=} Q(x) \cdot P_{12}(x) + R(x)$$

Only $R$ requires commitment, and only $F$ and $R$ need full 12C evaluation; the three sparse lines cost 4C each: \textbf{41C and 12 $\Fq$ commitments} --- a \textbf{46\% constraint reduction} and \textbf{75\% fewer commitments}. The comparison across all three schemes is summarised in Table~\ref{tab:miller_loop_comparison}.

\begin{table}[h]
\centering
\caption{3-pair Miller loop zero-bit step comparison}
\label{tab:miller_loop_comparison}
\begin{tabular}{|l|c|c|c|}
\hline
\textbf{Scheme} & \textbf{Constraints} & \textbf{$\Fq$ Commitments} & \textbf{vs baseline} \\
\hline
\cite{PR1CS22} & 132C & --- & --- \\
\hline
\cite{FXFM23} & 76C & 48 & $-$42\% \\
\hline
\textbf{Ours} & \textbf{41C} & \textbf{12} & $-$\textbf{69\%} \\
\hline
\multicolumn{4}{|l|}{\footnotesize \cite{PR1CS22}: $36\text{C} + 3\times30\text{C} = 132\text{C}$} \\
\multicolumn{4}{|l|}{\footnotesize \cite{FXFM23}: $5\times12\text{C} + 3\times4\text{C} + 4\text{C} = 76\text{C}$} \\
\multicolumn{4}{|l|}{\footnotesize Ours: $2\times12\text{C} + 3\times4\text{C} + 5\text{C} = 41\text{C}$\quad (excl.\ Fiat-Shamir overhead)} \\
\hline
\end{tabular}
\end{table}


\subsection{Operation Polynomials}\label{sec:op-polys}

We now state the verification equation for each step of Algorithm~\ref{alg:pairing} in a 3-pairing configuration (as used in Groth16 verification). All equations share the form $\text{LHS} = R(x) + Q(x) \cdot P_{12}(x)$; the quotient $Q$ is amortized across all steps via batching (Section~\ref{sec:iop:batching}).

\subsubsection{Zero-bit step}\label{sec:op:ml-0bit}

Lines~\ref{pr:ml:sq}--\ref{pr:ml:ldbl}: one squaring and three doubling lines.

\begin{equation}\label{eq:op:0b}
F^2(x) \cdot L_1(x) \cdot L_2(x) \cdot L_3(x) = R(x) + Q(x) \cdot P_{12}(x)
\end{equation}

$F, R \in \Fe{12}$ (12 coefficients each); $L_1, L_2, L_3 \in \text{Sparse}_{01379}$ (4 coefficients each). \textbf{Cost: 41C, 12 commitments.}

\subsubsection{Non-zero bit step}\label{sec:op:ml-nzbit}

Lines~\ref{pr:ml:rw1bit}--\ref{pr:ml:ladd-1bit}: squaring, residue witness, and both doubling and addition lines for each pair.

\begin{equation}\label{eq:op:n0b}
F^2(x) \cdot W^{\pm1}(x) \cdot L_{\text{d1}}(x) \cdot L_{\text{d2}}(x) \cdot L_{\text{d3}}(x) \cdot L_{\text{a1}}(x) \cdot L_{\text{a2}}(x) \cdot L_{\text{a3}}(x) = R(x) + Q(x) \cdot P_{12}(x)
\end{equation}

$W^{-1}$ is used for negative bits; the equation structure is identical. $W \in \Fe{12}$ encapsulates the equivalence class relationship from~\cite{OPP24}. \textbf{Cost: 69C, 12 commitments.}

\subsubsection{Frobenius correction step}\label{sec:op:ml-frob}

Lines~\ref{pr:ml:pi1}--\ref{pr:ml:pi2}: no squaring; two sparse correction lines per pair for the $\pi_p$ and $-\pi_p^2$ mappings.

\begin{equation}\label{eq:op:fm}
F(x) \cdot L_{1\pi}(x) \cdot L_{2\pi}(x) \cdot L_{3\pi}(x) \cdot L_{1\pi^2}(x) \cdot L_{2\pi^2}(x) \cdot L_{3\pi^2}(x) = R(x) + Q(x) \cdot P_{12}(x)
\end{equation}

\textbf{Cost: 55C, 12 commitments.}

\subsubsection{Final exponentiation step}\label{sec:op:finx}

Lines~\ref{pr:ml:27r}--\ref{pr:ml:fe}: cubic scaling and the Frobenius component ($q - q^2 + q^3$) of the final exponentiation, following~\cite{OPP24}. $R = 1$ for a valid pairing.

\begin{equation}\label{eq:op:fx}
F(x) \cdot c^{-q}(x) \cdot c^{q^2}(x) \cdot c^{-q^3}(x) \cdot W^w(x) = 1 + Q(x) \cdot P_{12}(x)
\end{equation}

$w \in \{0,1,2\}$ selects the cubic root of unity scaling; $W^w$ is sparse with 2 nonzero coefficients. The Frobenius images of $c$ are provided as prover inputs. \textbf{Cost: 56C, 12 commitments.}

\subsection{Security}\label{sec:op-sec}

Each verification equation takes the form $P = \prod_i A_i(x) - Q(x) \cdot P_{12}(x) - R(x)$ where $\deg(P) < 2^8$.\footnote{For $k$-pair pairings, $\deg(P) \approx 88 + 18(k-3)$, well below $2^8$ for all practical $k$. This bound can be extended to $2^{10}$ with negligible soundness impact.} By the Schwartz-Zippel lemma (Section~\ref{sec:bg:sz}):

\textbf{Soundness:} $\delta_s \leq d/|\Fq| < 2^8/2^{254} = 2^{-246}$ for BN254.

\textbf{Completeness:} An honest prover always produces $P(x) = 0$ identically, so the verifier always accepts. Euclidean division guarantees unique $Q, R$ with $\deg(R) < \deg(P_{12}) = 12$.